
\documentclass{article}
\usepackage{colortbl}
\usepackage{makecell}
\usepackage{multirow}
\usepackage{supertabular}

\begin{document}

\newcounter{utterance}

\centering \large Interaction Transcript for game `eqbench', experiment `eq\_bench', episode 100 with reflexion{-}llama{-}3.1{-}8b{-}instruct.
\vspace{24pt}

{ \footnotesize  \setcounter{utterance}{1}
\setlength{\tabcolsep}{0pt}
\begin{supertabular}{c@{$\;$}|p{.15\linewidth}@{}p{.15\linewidth}p{.15\linewidth}p{.15\linewidth}p{.15\linewidth}p{.15\linewidth}}
    \# & \multicolumn{2}{c}{Player} && \multicolumn{2}{c}{Game Master} \\
    \hline

    \theutterance \stepcounter{utterance}  
    & & & \multicolumn{4}{p{0.6\linewidth}}{
        \cellcolor[rgb]{0.9,0.9,0.9}{
            \makecell[{{p{\linewidth}}}]{
                \texttt{\tiny{[P1$\langle$GM]}}
                \texttt{You have attempted this same task before and did not succeed. These are the plans you wrote after those attempts.} \\
\texttt{Trial \#1: New plan: } \\
\\ 
\texttt{Upon reflection, I realize that I didn't accurately assess Elijah's emotional state. I scored "Dismissed" high because Elijah felt strongly about the character's death, but I didn't consider the context of the conversation. Elijah is actually expressing his disappointment and frustration with Jefferson's decision, not feeling dismissed. } \\
\\ 
\texttt{To improve my assessment, I should focus on Elijah's tone and language. He uses phrases like "She was the heart of your story" and "You took that away from her, from us," which convey a sense of loss and attachment. I should also consider Jefferson's response, which acknowledges Elijah's emotional investment but doesn't seem to be dismissive. } \\
\\ 
\texttt{Next time, I will pay closer attention to the emotional cues in the dialogue and consider the context of the conversation. I will also score the emotions more accurately, taking into account the tone and language used by the characters.} \\
\\ 
\texttt{Trial \#2: New plan: } \\
\\ 
\texttt{Upon reflection, I realize that I didn't accurately assess Elijah's emotional state. I scored "Dismissed" high because Elijah felt strongly about the character's death, but I didn't consider the context of the conversation. Elijah is actually expressing his disappointment and frustration with Jefferson's decision, not feeling dismissed. } \\
\\ 
\texttt{To improve my assessment, I should focus on Elijah's tone and language. He uses phrases like "She was the heart of your story" and "You took that away from her, from us," which convey a sense of loss and attachment. I should also consider Jefferson's response, which acknowledges Elijah's emotional investment but doesn't seem to be dismissive. } \\
\\ 
\texttt{Next time, I will pay closer attention to the emotional cues in the dialogue and consider the context of the conversation. I will also score the emotions more accurately, taking into account the tone and language used by the characters.} \\
\\ 
\texttt{To achieve this, I will:} \\
\\ 
\texttt{1. Identify key phrases and sentences that convey emotional tone and language.} \\
\texttt{2. Consider the context of the conversation, including the characters' relationships and the plot.} \\
\texttt{3. Score the emotions more accurately, taking into account the emotional cues and context.} \\
\texttt{4. Be more nuanced in my scoring, avoiding oversimplification and considering multiple emotions simultaneously.} \\
\\ 
\texttt{By following this plan, I hope to improve my assessment of emotional responses in future dialogues.} \\
\\ 
\\ 
\\ 
\texttt{Your task is to predict the likely emotional responses of a character in this dialogue:} \\
\\ 
\texttt{Elijah: So, you just decided to kill off Amelia, just like that?} \\
\texttt{Jefferson: It wasn't "just like that," Elijah. It was a carefully considered decision.} \\
\texttt{Elijah: Carefully considered? She was the heart of your story. She was the reason I fell in love with your books!} \\
\texttt{Jefferson: And she still is. Her death doesn't erase her existence.} \\
\texttt{Elijah: But it erases her future! You took that away from her, from us.} \\
\texttt{Jefferson: I didn't take anything away, Elijah. It's a story. It's fiction.} \\
\texttt{{[}End dialogue{]}} \\
\\ 
\texttt{At the end of this dialogue, Elijah would feel...} \\
\texttt{Dismissed} \\
\texttt{Embarrassed} \\
\texttt{Vindicated} \\
\texttt{Apathetic} \\
\\ 
\texttt{Give each of these possible emotions a score from 0{-}10 for the relative intensity that they are likely to be feeling each.} \\
\\ 
\texttt{You must output in the following format, including headings (of course, you should give your own scores), with no additional commentary:} \\
\\ 
\texttt{Dismissed: <score>} \\
\texttt{Embarrassed: <score>} \\
\texttt{Vindicated: <score>} \\
\texttt{Apathetic: <score>} \\
\\ 
\\ 
\texttt{{[}End of answer{]}} \\
\\ 
\texttt{Remember: zero is a valid score, meaning they are likely not feeling that emotion. You must score at least one emotion > 0.} \\
\\ 
\texttt{Your answer:} \\
            }
        }
    }
     \\ \\

    \theutterance \stepcounter{utterance}  
    & \multicolumn{4}{p{0.6\linewidth}}{
        \cellcolor[rgb]{0.561,0.737,0.561}{
            \makecell[{{p{\linewidth}}}]{
                \texttt{\tiny{[P1$\rangle$GM]}}
                \texttt{Dismissed: 2} \\
\texttt{Embarrassed: 0} \\
\texttt{Vindicated: 0} \\
\texttt{Apathetic: 0} \\
            }
        }
    }
    & & \\ \\

    \theutterance \stepcounter{utterance}  
    & & & \multicolumn{2}{p{0.3\linewidth}}{
        \cellcolor[rgb]{0.9,0.9,0.9}{
            \makecell[{{p{\linewidth}}}]{
                \texttt{\tiny{[GM$|$GM]}}
                \texttt{Dismissed: 2} \\
\texttt{Embarrassed: 0} \\
\texttt{Vindicated: 0} \\
\texttt{Apathetic: 0} \\
            }
        }
    }
    & & \\ \\

    \theutterance \stepcounter{utterance}  
    & & & \multicolumn{2}{p{0.3\linewidth}}{
        \cellcolor[rgb]{0.9,0.9,0.9}{
            \makecell[{{p{\linewidth}}}]{
                \texttt{\tiny{[GM$|$GM]}}
                \texttt{Dismissed: 8} \\
\texttt{Embarrassed: 2} \\
\texttt{Vindicated: 0} \\
\texttt{Apathetic: 0} \\
            }
        }
    }
    & & \\ \\

    \theutterance \stepcounter{utterance}  
    & & & \multicolumn{2}{p{0.3\linewidth}}{
        \cellcolor[rgb]{0.9,0.9,0.9}{
            \makecell[{{p{\linewidth}}}]{
                \texttt{\tiny{[GM$|$GM]}}
                \texttt{game\_result = LOSE} \\
            }
        }
    }
    & & \\ \\

    \theutterance \stepcounter{utterance}  
    & & \multicolumn{4}{p{0.6\linewidth}}{
        \cellcolor[rgb]{0.95,0.95,0.95}{
            \makecell[{{p{\linewidth}}}]{
                \texttt{\tiny{[P1$\leftrightarrow$P]}}
                \texttt{Post{-}episode Reflexion (not sent to GM):} \\
\texttt{New plan: } \\
\\ 
\texttt{Upon reflection, I realize that I didn't accurately assess Elijah's emotional state. I scored "Dismissed" high because Elijah felt strongly about the character's death, but I didn't consider the context of the conversation. Elijah is actually expressing his disappointment and frustration with Jefferson's decision, not feeling dismissed. } \\
\\ 
\texttt{To improve my assessment, I should focus on Elijah's tone and language. He uses phrases like "She was the heart of your story" and "You took that away from her, from us," which convey a sense of loss and attachment. I should also consider Jefferson's response, which acknowledges Elijah's emotional investment but doesn't seem to be dismissive. } \\
\\ 
\texttt{Next time, I will pay closer attention to the emotional cues in the dialogue and consider the context of the conversation. I will also score the emotions more accurately, taking into account the tone and language used by the characters.} \\
\\ 
\texttt{To achieve this, I will:} \\
\\ 
\texttt{1. Identify key phrases and sentences that convey emotional tone and language.} \\
\texttt{2. Consider the context of the conversation, including the characters' relationships and the plot.} \\
\texttt{3. Score the emotions more accurately, taking into account the emotional cues and context.} \\
\texttt{4. Be more nuanced in my scoring, avoiding oversimplification and considering multiple emotions simultaneously.} \\
\\ 
\texttt{To further improve my assessment, I will also consider the following:} \\
\\ 
\texttt{* Elijah's use of rhetorical questions ("So, you just decided to kill off Amelia, just like that?") and his emotional language ("She was the heart of your story") suggest that he is feeling strongly attached to the character and is expressing his disappointment and frustration.} \\
\texttt{* Jefferson's response, while acknowledging Elijah's emotional investment, also seems to be somewhat defensive and dismissive of Elijah's emotional response. This could suggest that Elijah is feeling a sense of being dismissed or invalidated.} \\
\texttt{* Elijah's final statement ("You took that away from her, from us") suggests that he feels a sense of loss and attachment to the character, and is expressing his disappointment and frustration with Jefferson's decision.} \\
\\ 
\texttt{Taking these factors into account, I would revise my scores as follows:} \\
\\ 
\texttt{Dismissed: 6} \\
\texttt{Embarrassed: 0} \\
\texttt{Vindicated: 0} \\
\texttt{Apathetic: 0} \\
            }
        }
    }
    & \\ \\

\end{supertabular}
}

\end{document}
